\begin{helper}[Transport of Tashkinov witnesses through a controlled component swap]
Let \(T=(y_0,e_1,y_1,\ldots,e_n,y_n)\) be a Tashkinov tree for a
proper partial coloring \(\varphi\).  Let \(i<j\), let
\(\alpha\ne\beta\) be missing at \(y_i,y_j\), respectively, and assume
that neither color occurs on \(e_2,\ldots,e_j\).  Swap the two colors on
the \((\alpha,\beta)\)-component \(C\) containing \(y_j\), obtaining
\(\varphi'\).  Then every edge \(e_2,\ldots,e_j\) keeps its exact color,
and every initial segment through \(y_j\) is a Tashkinov tree under
\(\varphi'\).  Moreover, if \(y_i\in C\), then every old witness for
every edge at or after position \(j\) gives a possibly exchanged earlier
witness after the swap, and the full displayed sequence \(T\) is a
Tashkinov tree under \(\varphi'\).
\end{helper}
\begin{proof}
Let \(h\) satisfy \(2\le h\le j\).  By the two
\noderef{TreeColorUnusedThrough} assumptions, the exact value
\(\varphi(e_h)\) is neither \(\alpha\) nor \(\beta\).  Whether or not
the edge-node belongs to \(C\), the definition
\noderef{TwoColorComponentSwap} therefore fixes that value.  This gives
the first assertion.

Let \(S\) be an initial segment through \(y_j\).  For any colored
displayed edge of \(S\), choose its earlier witness from
\noderef{TashkinovTree}.  The edge color is outside
\(\{\alpha,\beta\}\) by the preceding paragraph.  Its exact edge value
is fixed, and its missing status at the old witness is fixed by the last
clause of \noderef{TwoColorComponentSwapMissingBehavior}.  The vertex
and edge identities in \(S\) agree with those of \(T\) by
\noderef{TreeSequenceInitialSegment}.  Hence the same earlier vertex is
a valid witness in \(S\), proving that \(S\) is Tashkinov for
\(\varphi'\).

Assume now that \(y_i\in C\).  The root \(y_j\) is also in \(C\).
By \noderef{TwoColorComponentSwapMissingBehavior}, \(y_i\) misses
\(\beta\) and \(y_j\) misses \(\alpha\) after the swap.  Fix a later
edge \(e_h\), an old color \(c\), and an old earlier witness \(y_r\).
We construct the claimed new witness by exhaustive cases.

If \(c\notin\{\alpha,\beta\}\), the edge still has color \(c\), and
the same vertex still misses \(c\), by
\noderef{TwoColorComponentSwapMissingBehavior}.  If the edge-node lies
in \(C\) and its old color is \(\alpha\), its new color is \(\beta\)
and \(y_i\) is a new missing-\(\beta\) witness.  If its old color is
\(\beta\), the symmetric witness is \(y_j\).  Both vertices precede
\(e_h\), since \(i<j\le h\).

It remains that the edge-node is outside \(C\) and
\(c\in\{\alpha,\beta\}\), so its exact color is fixed.  If the old
witness is outside \(C\), its entire missing set is unchanged by
\noderef{TwoColorComponentSwapMissingBehavior}.  If it is in \(C\),
then it has degree at most one in the two-color incidence graph by
properness and \noderef{MissingColorIncidenceEndpoint}.  The same is
true of \(y_i\) and \(y_j\).  These two vertices are distinct because
\(i<j\) and \noderef{TreeSequence} makes the displayed vertices
injective.  By \noderef{TwoColorComponentSpanningPath}, a simple path
has support equal to all of \(C\).  A support vertex internal to that
path has two distinct path-neighbors, so each of the three
degree-at-most-one vertices must be an end.  Since the nontrivial path
has exactly two ends and \(y_i,y_j\) already occupy them, the old
witness must be one of \(y_i,y_j\).  Thus a lost missing-\(\alpha\)
witness can only be
\(y_i\), and it is replaced by the new missing-\(\alpha\) witness
\(y_j\).  Symmetrically, a lost missing-\(\beta\) witness at \(y_j\)
is replaced by \(y_i\).  Again both replacements occur before
\(e_h\).

This proves the quantified witness-transport assertion.  Applying it to
the witness supplied by \noderef{TashkinovTree} for each later edge, and
using the already proved initial-segment assertion for earlier edges,
shows that the full sequence \(T\) remains a Tashkinov tree.
\end{proof}
